\documentclass[a4paper,10pt,twoside]{article}
\usepackage[finnish]{babel}    % suomalaiset babel-makrot
\usepackage[utf8]{inputenc}  % Unixin merkist�
\usepackage[T1]{fontenc}       % kirjaimet, joissa aksentteja
\usepackage{lmodern}           % paremmat fontit
\usepackage{amsmath}           % AMS-paketteja
\usepackage{amsthm}
\usepackage{amsfonts}         % eiv�t toimi yhdess� mathdesignin kanssa
\usepackage{amssymb}          %
\usepackage{verbatim}          % comment-ymp�rist�j� varten
\usepackage{calc}              % tarvitaan mittojen v�lisiin laskutoimituksiin
\usepackage{graphicx}          % kuvia varten
\usepackage[all]{xy}           % kaavioita varten
%\usepackage{mparhack}          % korjaa marginaalihuomiot oikeaan (ei v��r��n) reunaan
%\usepackage{perpage}           % alaviitteet sivuittain
%\usepackage{chngcntr}          % voi lis�t� tai poistaa numeroinnin riippuvuuden toisesta
\usepackage{booktabs}          % paremmat taulukot
\usepackage{pbox}
\usepackage{eurosym}		%euro
\usepackage{textcomp} %punta


\addtolength{\textwidth}{2.5cm}      	% Suurennetaan sivun mittoja.
\addtolength{\hoffset}{-1.25cm}
\addtolength{\textheight}{5cm}
\addtolength{\voffset}{-2.5cm}
%\addtolength{\hoffset}{-1.0cm}
%\addtolength{\textwidth}{2.0cm}
\addtolength{\marginparsep}{0.2cm}
\setlength{\parskip}{0.2cm}


%%\usepackage{geometry}
%%\usepackage[hmarginratio=1:1]{geometry}
%%\geometry{bindingoffset=0.5cm}
%%\setlength{\oddsidemargin}{22pt}
%\setlength{\evensidemargin}{68pt}
%%\setlength{\oddsidemargin}{32pt}
%\setlength{\voffset}{-7mm}
%\setlength{\headsep}{1cm}
%\setlength{\footskip}{2.3cm}
%\setlength{\hoffset}{-5mm}
%\setlength{\textheight}{595pt}
%\setlength{\topmargin}{22pt}
%\setlength{\headheight}{14pt}
%%\setlength{\textwidth}{388pt}

% vain sectionin numero esiintyy luvun otsikossa (ei chapterin numeroa mukana)
% (t�hti ilmoittaa, ett� komento ei sis�ll� kappaleenvaihtoa)
\renewcommand*\thesection{\arabic{section}}

% vain yht�l�n numero esiintyy yht�l�rivill� (ei chapterin numeroa mukana)
\renewcommand*\theequation{\arabic{equation}}

% sectionien numerointi jatkuu, vaikka tulisi uusi chapterin numero
%\counterwithout{section}{chapter}

% kuvien numerointi alkaa alusta uudessa sectionissa
%\makeatletter
%\@addtoreset{figure}{section}
%\makeatother
%\counterwithin{figure}{section}
% sectionin numero esiintyy kuvien numeroinnissa
\renewcommand*\thefigure{\thesection.\arabic{figure}}
\renewcommand*\thetable{\thesection.\arabic{table}}

% alaviitteet sivuittain
%\MakePerPage{footnote}
% alaviitteet symboleiksi
%\renewcommand*\thefootnote{\fnsymbol{footnote}}


%\swapnumbers

\theoremstyle{plain}
\newtheorem{lause}{Lause}[subsection]
\newtheorem{lemma}[lause]{Lemma}
\newtheorem{prop}[lause]{Propositio}
\newtheorem{kor}[lause]{Korollaari}
%\newtheorem{lemmaIlmanNumeroa}{Lemma}
%\newtheorem{lauseIlmanNumeroa}{Lause}

\theoremstyle{definition}
\newtheorem{maar}[lause]{M\"a\"aritelm\"a}
\newtheorem{teht}[lause]{Teht�v�}
\newtheorem{esim}[lause]{Esimerkki}

\theoremstyle{remark}
\newtheorem*{lisa}{Lis�tieto}
\newtheorem*{huom}{Huomautus}
\newtheorem*{hist}{Historiaa}
\newtheorem*{ratk}{Ratkaisu}

% sijoitus integraaliin
\def\viiva #1 #2 {\mathop{\Big/}\limits_{\!\!\!{#1}}^{\>\,{#2}}}

\DeclareMathOperator{\kernel}{Ker}
\DeclareMathOperator{\image}{Im}
\DeclareMathOperator{\syt}{syt}
\DeclareMathOperator{\pyj}{pyj}
\DeclareMathOperator{\sgn}{sgn}
\providecommand{\clock}{\ensuremath{\oplus}}
\DeclareMathOperator{\proj}{\mathrm{proj}}
\DeclareMathOperator{\vir}{\mathrm{span}}
\DeclareMathOperator{\dist}{\mathrm{d}}
\DeclareMathOperator{\koht}{\mathrm{perp}}
\DeclareMathOperator{\kuvausL}{L}

\newcommand{\N}{\mathbb{N}}
\newcommand{\Z}{\mathbb{Z}}
\newcommand{\Q}{\mathbb{Q}}
\newcommand{\R}{\mathbb{R}}
\newcommand{\C}{\mathbb{C}}
\newcommand{\vP}{\mathcal{P}}
\newcommand{\kB}{\mathcal{B}}
\newcommand{\kC}{\mathcal{C}}
\newcommand{\kS}{\mathcal{S}}
\newcommand{\kR}{\mathcal{R}}
\newcommand{\kT}{\mathcal{T}}
\newcommand{\kE}{\mathcal{E}}
\newcommand{\teuro}{\text{ \euro}}
\newcommand{\tteuro}{\,\euro \ }
\newcommand{\vF}{\mathcal{F}}

\newcommand{\vx}{\bar{x}}
\newcommand{\vy}{\bar{y}}
\newcommand{\vz}{\bar{z}}
\newcommand{\vv}{\bar{v}}
\newcommand{\vw}{\bar{w}}
\newcommand{\vu}{\bar{u}}
\newcommand{\vp}{\bar{p}}
\newcommand{\vq}{\bar{q}}
\newcommand{\vn}{\bar{n}}
\newcommand{\va}{\bar{a}}
\newcommand{\vb}{\bar{b}}
\newcommand{\vi}{\bar{i}}
\newcommand{\vj}{\bar{j}}
\newcommand{\vk}{\bar{k}}
\newcommand{\ve}{\bar{e}}
\newcommand{\suoral}{\mathcal{l}}
\newcommand{\nv}{\bar{0}}
\newcommand{\id}{\mathrm{id}}
\newcommand{\pv}{\overline}
\newcommand{\tp}[1]{#1^{\top}}

%t�ydennetyt matriisit
\makeatletter
\renewcommand*\env@matrix[1][*\c@MaxMatrixCols c]{%
  \hskip -\arraycolsep
  \let\@ifnextchar\new@ifnextchar
  \array{#1}}
\makeatother

%%%%%%%%%%%%%%%%%%%%%%%%
%Fontteja:

%KP sans serif:
%\usepackage[sfmath]{kpfonts} 
%(Sfmath option only to make math in sans serif. Probablye only for use when base font is sans serif.)
%\renewcommand*\familydefault{\sfdefault} %K�ytet��n KP sans serifin kanssa 
%(Only if the base font of the document is to be sans serif)

%Kurier (condensed on hiukan "kapeampi" versio):
%\usepackage[condensed,math]{kurier} 
\usepackage[math]{kurier} 
%\usepackage{cmbright}

%Iwona:
%\usepackage[math]{iwona} 

%%%%%%%%%%%%%%%%%%%%%%%%

\newcommand{\gen}[1]{\langle #1 \rangle}
%\renewcommand{\bar}[1]{\overline{ #1 }}
\newcommand{\norm}[1]{\lVert #1 \rVert}

% alaviitteet symboleiksi
%\renewcommand{\thefootnote}{\fnsymbol{footnote}}

\newcommand{\marginaali}[1]{\marginpar[\raggedleft\textbf{#1}]{\raggedright\textbf{#1}}}
%\newcommand{\marginaali}[1]{\marginpar[\begin{center}\sffamily{#1}\end{center}]{\raggedright\textbf{#1}}}
%\newcommand{\marginaali}[1]{\marginpar[\raggedleft\sffamily\fontsize{9}{11}\selectfont{#1}]{\raggedright\sffamily\fontsize{9}{11}\selectfont{#1}}}
% \newcommand{\marginaali}[1]{\marginpar[\raggedleft\mbox{}{\tiny #1}]{\raggedright\mbox{}{\tiny #1}}}

%\newcommand{\maaritelma}[1]{
%  \noindent
%  \fbox{\parbox{0.97\textwidth}{
%      \begin{maar} #1 \end{maar}
%  }}
%  \vspace{0.3cm}
%}

\newcommand{\maaritelma}[2][\@empty]{
  \vspace{0.3cm}
  \noindent
  \marginaali{#1}
  \fbox{
    \parbox{0.95\textwidth}{
%      \vspace{-0.2cm}
%      \begin{maar} #2 \end{maar}
			#2
%      \vspace{-0.2cm}
    }
%    \begin{minipage}{0.95\textwidth}
%      \vspace{0.1cm}
%      \begin{maar} #2 \end{maar}
%      \vspace{-0.2cm}
%    \end{minipage}
  }
  \vspace{0.3cm}
}

\usepackage{float}
\usepackage{multicol}
\setlength{\multicolsep}{1ex}

\floatstyle{plain}
\newfloat{lisaosuus}{hbt}{lop}
\floatname{lisaosuus}{Lis�huomio}

\newcommand{\lisatieto}[1]{%
%\vspace{0.3cm}
\noindent
\begin{center}
\fbox{%
\begin{minipage}{\textwidth-0.6cm}
%\setlength{\parindent}{0.5cm}
\vspace{0.1cm}
\small
\begin{lisa}{#1}\vspace{0.1cm}\end{lisa}
\end{minipage}%
}
\end{center}
%\vspace{0.3cm}
}

\newenvironment{itemize*}{
  \begin{list}{$\bullet$}{
    \setlength{\itemsep}{0cm}
    \setlength{\topsep}{0cm}
    \setlength{\labelwidth}{3em}
    \setlength{\leftmargin}{\labelwidth+\labelsep}
  }
}
{\end{list}}

\let\olddescription\description
\renewcommand{\description}{
  \olddescription
  \setlength{\itemsep}{0pt}
  \setlength{\topsep}{0pt}
  \setlength{\itemindent}{0mm}
  %\setlength{\parskip}{0pt}
  %\setlength{\parsep}{0pt}
  %\setlength{\labelwidth}{1cm}
  %\setlength{\leftmargin}{2cm}
}

\newenvironment{viivoita*}{
  \begin{list}{--}{
    \setlength{\itemsep}{0cm}
    \setlength{\topsep}{0cm}
    \setlength{\labelwidth}{3em}
    \setlength{\leftmargin}{\labelwidth+\labelsep}
  }
}
{\end{list}}

\newenvironment{enumerate*}{
  \begin{enumerate}
  \setlength{\itemsep}{0cm}
  \setlength{\topsep}{0cm}
  \renewcommand{\theenumi}{\arabic{enumi}.}
  \renewcommand{\labelenumi}{\theenumi}
}
{\end{enumerate}}

\newenvironment{aakkosta*}{
  \begin{enumerate}
  \setlength{\itemsep}{0cm}
  \setlength{\topsep}{0cm}
  \setlength{\partopsep}{0mm}
  \renewcommand{\theenumi}{(\alph{enumi})}
  \renewcommand{\labelenumi}{\theenumi}
}
{\end{enumerate}}

\newcounter{tehtava}[section]
\renewcommand{\thetehtava}{\arabic{section}.\arabic{tehtava}}

\newenvironment{tehtavat}{
  \small
  \begin{list}{\thetehtava.}{
    \usecounter{tehtava}
    %\setlength{\itemsep}{0cm}
%\setlength{\parsep}{0cm}
    \setlength{\labelwidth}{3em}
    \setlength{\leftmargin}{\labelwidth+\labelsep}
  % sis�ist� itemize� tai aakkostamista varten
%  \renewcommand{\theenumii}{\alph{enumii})}
%  \renewcommand{\labelenumii}{\theenumii}
  }
%  \renewcommand{\labelenumi}{\arabic{enumi}.}
}
{\end{list}}

%%%%%%%%%%%%%%%%%%%%%%%%%%%%%%%%%%%%%%%%%%%%

% Alter some LaTeX defaults for better treatment of figures:
    % See p.105 of "TeX Unbound" for suggested values.
    % See pp. 199-200 of Lamport's "LaTeX" book for details.
    %   General parameters, for ALL pages:
    \renewcommand{\topfraction}{0.9}	% max fraction of floats at top
    \renewcommand{\bottomfraction}{0.8}	% max fraction of floats at bottom
    %   Parameters for TEXT pages (not float pages):
    \setcounter{topnumber}{2}
    \setcounter{bottomnumber}{2}
    \setcounter{totalnumber}{4}     % 2 may work better (oli 4)
    \setcounter{dbltopnumber}{2}    % for 2-column pages
    \renewcommand{\dbltopfraction}{0.9}	% fit big float above 2-col. text
    \renewcommand{\textfraction}{0.07}	% allow minimal text w. figs
    %   Parameters for FLOAT pages (not text pages):
    \renewcommand{\floatpagefraction}{0.7}	% require fuller float pages
	% N.B.: floatpagefraction MUST be less than topfraction !!
    \renewcommand{\dblfloatpagefraction}{0.7}	% require fuller float pages

	% remember to use [htp] or [htpb] for placement

%%%%%%%%%%%%%%%%%%%%%%%%%%%%%%%%%%%%%%%

\usepackage{multicol}
% Figures within a column...
\makeatletter
\newenvironment{tablehere}
{\def\@captype{table}}
{}
\newenvironment{figurehere}
{\def\@captype{figure}}
{}
\makeatother

%%%%%%%%%%%%%%%%%%%%%%%
\newcommand{\pa}{\mathopen]}		%avoin v�li
\newcommand{\pe}{\mathclose[}		%suljettu v�li
\newcommand{\isby}[1]{\stackrel{\textrm{\scriptsize #1}}{=} }
\newcommand{\funktio}[1]{\stackrel{\scriptstyle #1}{\to} }
\newcommand{\allev}[2]{\underline{\phantom{#2}}\llap{#1}}
\newcommand{\equivalent}{\ \Leftrightarrow\ }
\newcommand{\implikaatio}{\ \Rightarrow \ }

\setlength{\parindent}{0cm}
\renewcommand{\qedsymbol}{$\Box$}
%%%%%%%%%%%%%%%%%%%
\usepackage{fancyhdr}
\pagestyle{fancy}
\fancyhead{} % clear all header fields
\fancyhead[LE,RO]{\thepage}
\lfoot{}
\cfoot{\includegraphics[scale=0.35]{evlogo.eps}}
\rfoot{}
\renewcommand{\headrulewidth}{0pt}
\renewcommand{\footrulewidth}{0pt}
\usepackage{enumerate}

\newtheorem*{ratkaisu}{Ratkaisuehdotus}
\usepackage{environ}
\NewEnviron{killcontents}{}
\let\ratkaisu\killcontents
\let\endratkaisu\endkillcontents
\begin{document}
\pagestyle{empty}

\section*{Luentotehtäviä}

\subsection*{Korkolaskenta}

Huomioita
\begin{itemize}
    \item Lähdeveroa tai varsinkaan sen pyöristyssääntöä ei tarvitse ottaa tehtävissä huomioon, ellei toisin mainita.
    \item Jos yksinkertaisen korkolaskun tehtävissä on selvitettävä päivien lukumäärä, niin käytä yleistä laskutapaa todelliset/360, ellei toisin mainita.
\end{itemize}

\begin{enumerate}
\subsubsection*{Yksinkertainen korkolasku}

\item Mikä tilin korkokannan tulisi olla, jotta sinne talletettu 1000 euroa kasvaisi neljässä kuukaudessa korkoa 3,33 euroa?
\begin{ratkaisu}
    \[
    r = kit = 1000\cdot i\cdot \frac{4}{12} = 3,33\Leftrightarrow i = \frac{3,33}{1000\cdot\frac{4}{12}}\approx 0,01
    \]
    Siis korkokannan tulee olla 1~\%.
\end{ratkaisu}

\item Rafaelin tilin korkokanta oli kiinteä 1,3~\%. Hänen tallettamansa summa kasvoi yhdeksässä kuukaudessa korkoa 1,95 \euro.
Kuinka suuren summan hän oli alunperin tallettanut tililleen?
\begin{ratkaisu}
    \[
    r = kit \Leftrightarrow k = \frac{r}{it} = \frac{1,95}{0,013\cdot\frac{9}{12}} = 200
    \]
    Alkuperäinen pääoma on 200 euroa.
\end{ratkaisu}

\item Henkilö ottaa pakon edessä 500~\euro\ lainaa korkokannan ollessa 15~\% p.a. Kuinka pian velan suuruus on 550 euroa?
\begin{ratkaisu}
    \[
        r = kit \Leftrightarrow t = \frac{r}{ki} = \frac{50}{500\cdot0,15} = \frac{2}{3}
    \]
    Korkoajan on oltava 2/3-vuotta, eli 8 kuukautta.
\end{ratkaisu}

\item Ronja talletti 314 euroa tammikuun alussa tililleen. Hän nosti rahat kesäkuun lopussa, jolloin nettokorkoa oli kertynyt
 1,70 euroa. Mikä oli tilin nettokorkokanta? Entä korkokanta?
\begin{ratkaisu}
    Korkoa maksetaan kuudelta kuukaudelta. Nyt
    \[
        i = \frac{r}{kt} = \frac{1,70}{314\cdot\frac{6}{12}}\approx 0,01083
    \]
    joten nettokorkokanta  on likimain 1,1~\%-. Korkokanta on \(\frac{0,01083}{0,7}\approx 1,5~\%\).
\end{ratkaisu}

\item Anja nosti talletuksensa vuoden kuluttua talletushetkestä. Hän sai pankista verojen vähentämisen jälkeen 1073,20 euroa.
 Kuinka suuren summan hän oli tililleen tallettanut, kun tilin korkokanta oli 1,1~\% ja lähdevero 30~\%?
\begin{ratkaisu}
    \[
        K = k + r = k + kit = k(1+it)
    \]
    Nyt \(t = 1\), joten
    \[
        1073,20 = k(1+0,7\cdot0,011)\Leftrightarrow k = \frac{1073,20}{1+0,7\cdot0,011} \approx 1064,9995
    \]
    Alkuperäinen pääoma on siis 1065,00 euroa.
    \end{ratkaisu}

\item Saat yllättäen 75~000 euroa ja talletat koko summan 3.7.2014 pankkitilille, jonka korkokanta on 2{,}8~\%.
Nostat rahat tililtä 19.11.2014. Kuinka paljon saat nettokorkoa?
\begin{ratkaisu}
    Korkopäiviä kertyy \(31-3+31+30+31+19 = 139\). Nettokorko on
    \[
        r = kit = 75000\cdot0,7*0,028\cdot\frac{139}{360}\approx 567,58
    \]
    euroa.
\end{ratkaisu}

\item Kaverisi säästää rahaa pitkää lomareissua varten tallettamalla joka kuukauden alussa yhtä suuren summan rahaa tililleen.
Kaverisi aloittaa säästämisen tammikuun alussa ja haluaa nostaa rahat vuoden kuluttua.
Kuinka paljon rahaa hänen tulisi tallettaa tililleen kuukausittain, jotta vuoden kuluttua nostettavissa olisi 1 200 euroa?
Tilin vuosikorko on 2{,}5~\%. Muista lähdevero 30 \%.
\begin{ratkaisu}
    Olkoon \(k\) kuukausittain talletettava pääoma. Tammikuun talletus kasvaa korkoa 12 kk, helmikuun talletus 11 kk jne.
    \begin{align*}
        r &= k\cdot0,7\cdot0,025\cdot\left(\frac{12}{12} + \frac{11}{12} + \ldots + \frac{2}{12} + \frac{1}{12}\right)\\
         =& k\cdot0,7\cdot0,025\cdot\left(12\cdot\frac{\frac{12}{12} + \frac{1}{12}}{2}\right)\\
         =& k\cdot0,7\cdot0,025\cdot\frac{78}{12}\\
         =& 0,11375\cdot k
    \end{align*}
    Nyt \(1200 = 12k + r = 12k + 0,11375k = 12,11375k\) eli \(k = 1200/12,11375\approx 99,061\).
    Tilille on siis talletettava vähintään 99,07 euroa.
\end{ratkaisu}

\subsubsection*{Koronkorko}

\item Sijoitat 500 euroa sijoitusvakuutukseen, joka lupaa vähintään 3{,}0 \% vuotuisen koron sijoitukselle.
Neljän vuoden kuluttua sijoituksen arvo oli kasvanut 60 euroa. Toteutuiko sijoitusvakuutuksen lupaama vuosikorko?
\begin{ratkaisu}
    Lupauksen mukaan tililtä olisi nostettavissa \(500\cdot1,03^4\approx562,75\) euroa. Lupaus ei siis toteutunut.
\end{ratkaisu}

\item Viivi tallettaa saamansa 1800 euron palkan tilille, jonka korkokanta 2,1 \%. Kuinka suuren summan Viivi voi nostaa tililtä
neljän vuoden kuluttua? Muista lähdevero.
\begin{ratkaisu}
    Nettokorkokanta on \(0,7\cdot0,021 = 0,0147\),
    joten neljän vuoden kuluttua summa on
    \[
        1800\cdot1,0147^4\approx1908,1967\dots
    \]
    euroa. Tililtä on siis nostettavissa 1908,19 euroa.
\end{ratkaisu}

\item Mikä tulisi Viivin tilin korkokannan olla edellisen tehtävän tilanteessa, jotta tililtä olisi nostettavissa 2100 euroa?
\begin{ratkaisu}
    \[
        2100 = 1800\cdot q^4\Leftrightarrow q^4 = \frac{2100}{1800}
    \]
    Tästä saadaan \(q = 1,0393\dots\), joten nettokorkokannan tulisi olla 3,9~\% ja korkokannan 5,6~\%.
\end{ratkaisu}

\item Kaukainen sukulainen pyytää sinua avaamaan tilin, jotta hän voi jokaisen vuoden lopussa lahjoittaa sinulle tietyn
summan rahaa tallettamalla sen tälle tilille. Tilillä ei ole näiden vuosittaisten talletusten ja koronmaksun lisäksi muita tilitapahtumia.
Ensimmäisen talletuksen sukulaisesi tekee vuoden 2014 lopussa. Viiden vuoden kuluttua vuoden 2020 alussa tilillä on rahaa 4 342{,}15 euroa.
Kuinka suuri on vuotuinen talletus, jos tilin korko on 1{,}9 \% p.a? Huomioi lähdevero.
\begin{ratkaisu}
Olkoon \(x\) vuosittain talletettava (lahjoitettava) summa. Vuonna 2014 saatu summa kasvaa korkoa 5 kertaa,
vuonna 2015 saatu summa 4 kertaa jne., vuonna 2019 saatu summa ei ehdi kasvaa korkoa kertaakaan
(sillä talletus tehdään vuoden lopussa). Korkoa kasvaneiden talletusten summa on
\begin{align*}
    &x\cdot q^5 + x\cdot q^4+ x\cdot q^3 + x\cdot q^2+ x\cdot q + x \\
    =& x + x\cdot q+ x\cdot q^2 + x\cdot q^3+ x\cdot q^4 + x\cdot q^5\\
    =& x\cdot\frac{1-q^6}{1-q}
\end{align*}
joten
\[
   x\cdot\frac{1-q^6}{1-q} = 4324,15
\]
Tästä saadaan
\[
    x = \frac{4342,15\cdot(1-q)}{1-q^6}
\]
Sijoitetaan \(q= 1 + 0,7\cdot0,019=1,0133\) ja saadaan
\[
x = \frac{4342,15\cdot(1-1,0133)}{1-1,0133^6} \approx 699,99979
\]
Vuosittainen talletus on siis 700,00 euroa.
\end{ratkaisu}

\item Uuden vuoden lupauksena aloitat säästämisen ja talletat tammikuun lopusta alkaen joka toisen kuukauden lopussa 300 euroa tilille,
jonka korkokanta on 2{,}1 \% p.a. Kuinka paljon rahaa tililläsi on neljännen talletusvuoden lopussa koronmaksun jälkeen? Huomioi lähdevero.
\begin{ratkaisu}
Tarkastellaan ensin yhden vuoden talletuksia. Tammikuun talletus kasvaa korkoa 11 kuukautta, maaliskuun talletus 9 kuukautta,
toukokuun talletus 7 kuukautta jne., marraskuun talletus kasvaa korkoa yhden kuukauden. Vuoden talletuksille korkoa maksetaan yhteensä
\begin{align*}
    &300\cdot0,7\cdot0,021\cdot\left(\frac{11}{12} + \frac{9}{12} + \frac{7}{12} + \frac{5}{12} + \frac{3}{12} + \frac{1}{12}\right)\\
    =&300\cdot0,7\cdot0,021\cdot\frac{1}{12}\cdot(11+9+7+5+3+1)\\
    =& 13,23
\end{align*}
euroa. Tilannetta voidaan ajatella siis siten, että kunkin vuoden \emph{lopussa} tilille talletetaan \(6\cdot300+13,23 = 1813,23\) euroa.
Kaiken kaikkiaan tilillä on neljännen talletusvuoden lopussa rahaa
\[
    K = 1813,13\cdot q^3 + 1813,13\cdot q^2 + 1813,13\cdot q + 1813,13 = 1813,13\cdot\frac{1-q^4}{1-q}
\]
euroa. Sijoitetaan \(q = 1 + 0,7\cdot0,021 = 1,0147\) ja saadaan
\[
    K = 1813,13\cdot\frac{1-1,0147^4}{1-1,0147} \approx 7414,011022
\]
Tililtä on siis nostettavissa 7414,01 euroa.
\end{ratkaisu}

\subsubsection*{Diskonttaus}

\item Josefiina tallettaa tililleen vuoden alussa summan, jonka toivoo kasvavan ainakin 5000 euroksi kymmenessä vuodessa. Tilin korkokanta on 1,5~\%.
\begin{enumerate}[(a)]
    \item Selvitä diskonttaustekijä.
    \item Kuinka suuri summa Josefiinan tulee tililleen tallettaa?
\end{enumerate}
\begin{ratkaisu}
    \begin{enumerate}[(a)]
        \item Koska \(q = 1 + 0,0157 = 1,0157\), diskonttaustekijä on \(q^{-10} = \frac{1}{q^{10}}\approx 0,85574\).
        \item \(k = K\cdot q^{-10} = K\cdot 0,855747 = 5000\cdot 0,85574 \approx 4278,735\) euroa.
        Josefiinan tulee tallettaa tililleen vähintään 4278,74 euroa.
    \end{enumerate}
\end{ratkaisu}

\item Kauppias tarjoaa auton maksamiseen seuraavaa rahoitusmallia: käsiraha 8~000~\euro \ on maksettava heti ja tämän
jälkeen kahtena seuraavana vuonna maksetaan 5 000~\euro. Kolmen vuoden kuluttua maksetaan vielä 2 000~\euro.
Jos asiakas kuitenkin haluaa maksaa auton käteisellä heti ostohetkellä, mikä auton hinnan pitäisi olla, että asiakas
maksaisi siitä yhtä paljon kuin kauppiaan ehdottamassa rahoitusmallissa? Oletetaan, että kauppiaan käyttämä
korkokanta on 7{,}5 \%.
\begin{ratkaisu}
    Diskontataan maksuerät nykyhetkeen. Kerroin on \(q= 1+0,075 = 1,075\), joten nykyarvoksi saadaan
    \[
        8000 + \frac{5000}{1,075} + \frac{5000}{1,075^2} + \frac{2000}{1,075^3} = 18587,74699
    \]
    euroa. Auton hinta ostohetkellä tulisi olla siis 18587,75 euroa.
\end{ratkaisu}


\item Kuinka paljon lukion stipendirahastoon on lahjoitettava rahaa eurona, kun tarkoituksena on jakaa lahjoitus korkoineen
stipendeinä seuraavasti: tasan vuoden kuluttua lahjoituksesta 200 euroa, kahden vuoden kuluttua 300 euroa,
kolmen vuoden 400 euroa, neljän 500 euroa ja viiden vuoden kuluttua 600 euroa. Pääomaan lisätään vuosittain 4,5 prosentin korko,
ensimmäisen kerran vuoden kuluttua lahjoituksesta. [S2001, 14]
\begin{ratkaisu}
    Diskontataan maksuerät nykyhetkeen. Kerroin on \(q= 1+0,045 = 1,045\), joten nykyarvoksi saadaan
    \[
        \frac{200}{1,045} + \frac{300}{1,045^2} + \frac{400}{1,045^3} + \frac{500}{1,045^4} + \frac{600}{1,045^5} \approx 1717,38
    \]
    euroa.
\end{ratkaisu}

\item Ostat kalliin vekottimen, josta maksat kaupantekohetkellä 400 euron käsirahan ja kuukauden kuluttua 800 euroa korkokannan ollessa 6{,}6 \% p.a.
Kaverisi osti samanlaisen vekottimen, maksoi ostohetkellä 300 euroa ja neljän kuukauden kuluttua loput.
Kuinka paljon kaverisi maksoi toisessa maksuerässä, jos vekottimenne olivat yhtä kalliit?
\begin{ratkaisu}
    Oman hinnan nykyarvo on
    \[
        400 + \frac{800}{1+0,066\cdot\frac{1}{12}} = 400 + \frac{800}{1,0055} = 1195,624
    \]
    euroa. Kaverin hinnan nykyarvo
    \[
        300 + \frac{x}{1+0,066\cdot\frac{4}{12}} = 300 + \frac{x}{1,022}
    \]
    Jos hinnat ovat samat, niin
    \[
        300 + \frac{x}{1,022} =  1195,624
    \]
    eli \( x = 1,022\cdot(1195,624-300)\approx 915,32\) euroa.
\end{ratkaisu}

\item Tuttavasi on myymässä asuntoaan ja saa kaksi ostotarjousta. Ensimmäisessä tarjouksessa luvataan maksaa heti 145 000 euroa.
Toisessa tarjouksessa luvataan maksaa heti 64 000 euroa, vuoden kuluttua 54 000 euroa ja kahden vuoden kuluttua 30 000 euroa.
Kumpi tarjous on tuttavasi kannalta parempi, jos korkokanta on 4 \%? Entä, jos korkokanta on 1~\%?
\begin{ratkaisu}
    Diskontataan toisen tarjouksen summat nykyhetkeen. Jos korkokanta on 4~\%, niin nykyarvo on
    \[
        64000 + \frac{54000}{1,004} + \frac{30000}{1,004^2} \approx 143659,76
    \] euroa.
    Jos korkokanta on 1~\%, niin nykyarvo on vastaavasti
    \[
        64000 + \frac{54000}{1,001} + \frac{30000}{1,001^2} \approx 146874,23
    \] euroa.
    Näin ollen 4~\% korkokannalla kertamaksu on tuttavalle parempi, 1~\% korkokannalla osamaksu.
\end{ratkaisu}

\item Yritys on hankkimassa uuden tuotantohärvelin. Kauppias tarjoaa kaksi maksuvaihtoehtoa:
koko kauppahinta 15 000 euroa kaupanteon yhteydessä tai 7 000 euroa kaupanteon yhteydessä,
5 000 euroa vuoden kuluttua ja 3 500 euroa kahden vuoden kuluttua. Jos yritys maksaa investoinnin lainalla,
jonka vuotuinen korkokanta on 4{,}5 \%, kumpi maksuvaihtoehto sen kannattaa valita?
\begin{ratkaisu}
    Diskontataan toisen tarjouksen summat nykyhetkeen. Korkokanta on 4,5~\%, joten erien yhteenlaskettu nykyarvo on
    \[
        7000 + \frac{5000}{1,0045} + \frac{3500}{1,0045^2} \approx 14989,74
    \]
    euroa. Osamaksu tulee siis yritykselle halvemmaksi.
\end{ratkaisu}

\end{enumerate}

\end{document}